\documentclass[12pt, a4paper]{article}

% ── Shared preamble (src/lecturenotes.sty) ────────────────────────────────────
%   Packages, colors, boxes, \sol/\soleq and math macros all live in
%   src/*.sty. \input@path points LaTeX two folders up to find them.
\makeatletter
\def\input@path{{../../}}
\makeatother
\usepackage{lecturenotes}

% ── Local macros ──────────────────────────────────────────────────────────────
\providecommand{\bZ}{\mathbf{Z}}
\providecommand{\bH}{\mathbf{H}}

% ── Local TikZ styles (match S1, S2) ──────────────────────────────────────────
\usetikzlibrary{decorations.pathreplacing}
\tikzset{
  layerlabel/.style={font=\small\itshape, gray},
  opbox/.style={draw=orange!70!black, fill=notebg, thick, rounded corners=4pt,
      minimum height=0.9cm, minimum width=1.9cm, font=\small, align=center},
  window/.style={draw=orange!80!black, line width=1.6pt},
  brace/.style={decorate, decoration={brace, amplitude=4pt}, thick, gray},
}

% \drawgrid{<shift>}{<fill>}{<rows>} — grid of unit cells with values
\newcommand{\drawgrid}[3]{%
  \begin{scope}[shift={#1}]
    \foreach \row [count=\r] in {#3} {
        \foreach \v [count=\c] in \row {
            \draw[gray!70, fill=#2] (\c-1,-\r+1) rectangle ++(1,-1);
            \node[font=\small] at (\c-0.5,-\r+0.5) {$\v$};
          }
      }
  \end{scope}%
}

% \cuboid{<x>}{<depth: channels>}{<height: spatial>}{<color>}{<label>}
%   pseudo-3D block for a (C, H, W) tensor, centred vertically on y = 0
\newcommand{\cuboid}[5]{%
  \pgfmathsetmacro{\off}{#3*0.3}
  \filldraw[fill=#4!20, draw=#4!80!black] (#1,-#3/2) rectangle ++(#2,#3);
  \filldraw[fill=#4!35, draw=#4!80!black]
  (#1,#3/2) -- ++(\off,\off) -- ++(#2,0) -- ++(-\off,-\off) -- cycle;
  \filldraw[fill=#4!50, draw=#4!80!black]
  (#1+#2,-#3/2) -- ++(\off,\off) -- ++(0,#3) -- ++(-\off,-\off) -- cycle;
  \node[below, font=\scriptsize, align=center] at (#1+#2/2,-1.75) {#5};
}

% ── Document ──────────────────────────────────────────────────────────────────
\begin{document}

\begin{center}
  {\LARGE \textbf{Section 4.3: Building a CNN}} \\[0.3em]
  {\large\itshape Activation, Pooling and the Classifier Head} \\[0.5em]
  {\large AI-IM Class --- Lecture Notes}%
  \ifstudentmode
    \quad {\normalsize\color{gray}[Student Copy]}%
  \fi
\end{center}

\vspace{1em}
\hrule
\vspace{1.5em}

Section 4.2 gave us the convolutional layer. A real CNN wraps every
convolution with two more steps and finishes with the MLP classifier of
Section 3.4. In this section we
\begin{itemize}[nosep]
  \item describe the three stages of a typical CNN layer,
  \item compute max and average pooling by hand and see why pooling gives
        invariance to small shifts,
  \item assemble a complete CNN for MNIST and trace its shapes and
        parameters, and
  \item see how gradients pass back through pooling.
\end{itemize}

% ══════════════════════════════════════════════════════════════════════════════
\section{The Three Stages of a CNN Layer}

Goodfellow et al.\ describe a typical CNN layer as three stages applied in
sequence:

\begin{figure}[H]
  \centering
  \begin{tikzpicture}[>=Stealth, node distance=0.9cm]
    \node (in) {$\bX$};
    \node[opbox, right=of in] (conv) {convolution\\[-2pt]\scriptsize (affine)};
    \node[opbox, right=of conv] (relu) {detector\\[-2pt]\scriptsize (ReLU)};
    \node[opbox, right=of relu] (pool) {pooling\\[-2pt]\scriptsize (max)};
    \node[right=of pool] (out) {next layer};
    \draw[->, thick] (in) -- (conv);
    \draw[->, thick] (conv) -- node[above, font=\small] {$\bZ$} (relu);
    \draw[->, thick] (relu) -- node[above, font=\small] {$\bH$} (pool);
    \draw[->, thick] (pool) -- (out);
    \node[font=\scriptsize\ttfamily, gray] at (conv.south) [below=3pt] {nn.Conv2d};
    \node[font=\scriptsize\ttfamily, gray] at (relu.south) [below=3pt] {nn.ReLU};
    \node[font=\scriptsize\ttfamily, gray] at (pool.south) [below=3pt] {nn.MaxPool2d};
  \end{tikzpicture}
  \caption{The three stages of a CNN layer, with the matching PyTorch modules.}
  \label{fig:stages}
\end{figure}

\begin{enumerate}
  \item \textbf{Convolution} computes the pre-activations $\bZ$, as in
        Section 4.2. Like $\bW\bh + \bb$ in an MLP, it is linear (affine).
  \item \textbf{Detector stage.} A nonlinear activation, almost always ReLU,
        is applied to each entry: $\bH = \relu(\bZ)$. Without it, stacked
        convolutions would collapse into one linear map, exactly as stacked
        linear layers do in an MLP.
  \item \textbf{Pooling} replaces each small region of $\bH$ by a summary
        statistic, such as its maximum. This is new, and is the topic of the
        next section.
\end{enumerate}

\begin{stepbox}{Exercise 1 \quad The detector stage}
  A convolution produced the feature map $\bZ$ below. Apply ReLU.
  \[
    \bZ =
    \begin{pmatrix}
      1 & -2 & 3 & 0 \\ -1 & 4 & -3 & 2 \\ 0 & 1 & 5 & -4 \\ 2 & -1 & 1 & 3
    \end{pmatrix}
    \qquad\Longrightarrow\qquad
    \bH = \relu(\bZ) =
    \begin{pmatrix}
      \sol[0.5cm]{1} & \sol[0.5cm]{0} & \sol[0.5cm]{3} & \sol[0.5cm]{0} \\[2pt]
      \sol[0.5cm]{0} & \sol[0.5cm]{4} & \sol[0.5cm]{0} & \sol[0.5cm]{2} \\[2pt]
      \sol[0.5cm]{0} & \sol[0.5cm]{1} & \sol[0.5cm]{5} & \sol[0.5cm]{0} \\[2pt]
      \sol[0.5cm]{2} & \sol[0.5cm]{0} & \sol[0.5cm]{1} & \sol[0.5cm]{3}
    \end{pmatrix}
  \]
\end{stepbox}

% ══════════════════════════════════════════════════════════════════════════════
\section{Pooling}

A \textbf{pooling} layer slides a window over each feature map, just like a
convolution, but instead of a weighted sum it computes a fixed summary of the
window:
\begin{itemize}[nosep]
  \item \textbf{max pooling} takes the largest value (most common);
  \item \textbf{average pooling} takes the mean.
\end{itemize}
The usual choice is a $2 \times 2$ window with stride $2$, so the windows do
not overlap and each feature map shrinks to half its height and width.

\begin{figure}[H]
  \centering
  \begin{tikzpicture}[x=0.75cm, y=0.75cm, >=Stealth]
    \fill[orange!25]       (0,0)   rectangle (2,-2);
    \fill[mainblue!20]     (2,0)   rectangle (4,-2);
    \fill[green!25]        (0,-2)  rectangle (2,-4);
    \fill[purple!20]       (2,-2)  rectangle (4,-4);
    \foreach \row [count=\r] in {{1,0,3,0},{0,4,0,2},{0,1,5,0},{2,0,1,3}} {
        \foreach \v [count=\c] in \row {
            \draw[gray!70] (\c-1,-\r+1) rectangle ++(1,-1);
            \node[font=\small] at (\c-0.5,-\r+0.5) {$\v$};
          }
      }
    \draw[thick] (0,0) rectangle (4,-4);
    \draw[thick] (2,0) -- (2,-4); \draw[thick] (0,-2) -- (4,-2);
    \node[layerlabel] at (2,-4.7) {$\bH$ ($4\times4$)};

    \draw[->, thick] (4.8,-2) -- node[above, font=\small, align=center]
    {$2\times2$ max\\stride 2} (8.2,-2);

    \fill[orange!25]   (9,-1)  rectangle (10,-2);
    \fill[mainblue!20] (10,-1) rectangle (11,-2);
    \fill[green!25]    (9,-2)  rectangle (10,-3);
    \fill[purple!20]   (10,-2) rectangle (11,-3);
    \draw[thick] (9,-1) grid (11,-3);
    \node[layerlabel] at (10,-3.7) {output ($2\times2$)};
  \end{tikzpicture}
  \caption{Max pooling with a $2 \times 2$ window and stride $2$. Each
    coloured block of four values becomes one output value.}
  \label{fig:maxpool}
\end{figure}

\begin{stepbox}{Exercise 2 \quad Max and average pooling}
  Pool the feature map $\bH$ from Exercise 1 (Figure~\ref{fig:maxpool}) with
  a $2 \times 2$ window and stride $2$.
  \[
    \text{max pooling:}\;
    \begin{pmatrix}
      \sol[0.7cm]{4} & \sol[0.7cm]{3} \\[2pt]
      \sol[0.7cm]{2} & \sol[0.7cm]{5}
    \end{pmatrix},
    \qquad
    \text{average pooling:}\;
    \begin{pmatrix}
      \sol[0.9cm]{1.25} & \sol[0.9cm]{1.25} \\[2pt]
      \sol[0.9cm]{0.75} & \sol[0.9cm]{2.25}
    \end{pmatrix}
  \]
  What fraction of the values is discarded? \quad
  $\sol[3cm]{12/16 = 75\%}$
\end{stepbox}

Some facts about pooling layers:
\begin{itemize}
  \item A pooling layer has \textbf{no parameters}. Nothing is learned.
  \item It acts on \textbf{each channel separately}: a $(C, H, W)$ input gives
        a $(C, H', W')$ output with the same number of channels.
  \item The output size follows the same formula as convolution,
        $n_{\text{out}} = \lfloor (n + 2p - k)/s \rfloor + 1$. For $k = s = 2$
        and $p = 0$ this is $\lfloor n/2 \rfloor$.
  \item By shrinking the feature maps it cuts the computation of all later
        layers, and it makes the receptive field grow much faster than
        stacking convolutions alone (Section 4.1).
\end{itemize}

% ──────────────────────────────────────────────────────────────────────────────
\subsection{Pooling Gives Invariance to Small Shifts}

In Section 4.1 we saw that convolution is \emph{equivariant}: shift the
input and the feature map shifts with it. Pooling makes the representation
approximately \emph{invariant}: if the input moves a little, most pooled
values do not change at all. The reason is simple. The maximum of a window does
not care \emph{where} in the window the largest value is.

\begin{stepbox}{Exercise 3 \quad Shift the input by one pixel}
  A detector stage produced the 1-D output $\bh_A$. In $\bh_B$ the input was
  shifted one pixel to the right, so every value moved one place. Apply max
  pooling with window $3$ and stride $1$.
  \begin{center}
    \renewcommand{\arraystretch}{1.4}
    \begin{tabular}{l*{7}{c}}
      \toprule
      $\bh_A$         & 1 & 5 & 2 & 0 & 4 & 1 & 0 \\
      max pool        & $\sol[0.6cm]{5}$ & $\sol[0.6cm]{5}$ & $\sol[0.6cm]{4}$
                      & $\sol[0.6cm]{4}$ & $\sol[0.6cm]{4}$ & & \\
      \midrule
      $\bh_B$         & 0 & 1 & 5 & 2 & 0 & 4 & 1 \\
      max pool        & $\sol[0.6cm]{5}$ & $\sol[0.6cm]{5}$ & $\sol[0.6cm]{5}$
                      & $\sol[0.6cm]{4}$ & $\sol[0.6cm]{4}$ & & \\
      \bottomrule
    \end{tabular}
  \end{center}
  Between $\bh_A$ and $\bh_B$, $\sol[1cm]{7}$ of the $7$ detector values
  changed, but only $\sol[1cm]{1}$ of the $5$ pooled values changed.
\end{stepbox}

This invariance is \emph{local}: it holds for shifts smaller than the pooling
window. Stacking several conv--pool layers extends it to larger shifts, which
is what a classifier needs: the answer ``digit 7'' should not depend on
exactly where the 7 is written.

\begin{notebox}
  \textbf{Pooling over channels.} Pooling need not be over space. If several
  filters each detect the same shape at a different rotation, taking the
  maximum \emph{across those channels} gives a unit that responds to the
  shape at any of those rotations. The network can learn invariances that
  convolution alone does not provide (Goodfellow et al., Figure 9.9).
\end{notebox}

\begin{notebox}
  \textbf{Pooling or stride?} A convolution with stride $2$ also halves the
  feature map, and unlike pooling it has learned weights. Many modern
  architectures (e.g.\ ResNet, Section 4.4) shrink the maps mostly with
  strided convolutions and use pooling only at the start and end.
\end{notebox}

% ══════════════════════════════════════════════════════════════════════════════
\section{Assembling a Complete CNN}

A CNN for classification has two parts:
\begin{enumerate}
  \item a \textbf{feature extractor}: a stack of conv--ReLU--pool blocks.
        Each block halves the height and width and usually doubles the number
        of channels. The network trades spatial detail (\emph{where}) for
        richer features (\emph{what});
  \item a \textbf{classifier head}: flatten the final feature maps into a
        vector and feed it to the MLP classifier of Section 3.4, ending in
        $C$ logits, softmax and the cross-entropy loss.
\end{enumerate}

\begin{figure}[H]
  \centering
  \begin{tikzpicture}[>=Stealth, scale=0.95, every node/.style={transform shape}]
    \cuboid{0}{0.1}{2.8}{gray}{input\\$(1,28,28)$}
    \cuboid{1.9}{0.5}{2.8}{mainblue}{conv1\\$(16,28,28)$}
    \cuboid{4.1}{0.5}{1.4}{mainblue}{pool1\\$(16,14,14)$}
    \cuboid{5.9}{1.0}{1.4}{orange}{conv2\\$(32,14,14)$}
    \cuboid{8.3}{1.0}{0.7}{orange}{pool2\\$(32,7,7)$}

    \draw[->, thick] (1.0,0) -- (1.8,0);
    \draw[->, thick] (3.4,0) -- (4.0,0);
    \draw[->, thick] (5.1,0) -- (5.8,0);
    \draw[->, thick] (7.5,0) -- (8.2,0);

    % classifier head
    \filldraw[fill=green!20, draw=green!50!black] (10.4,-1.4) rectangle ++(0.25,2.8);
    \node[below, font=\scriptsize, align=center] at (10.52,-1.75) {flatten\\$1568$};
    \filldraw[fill=green!20, draw=green!50!black] (11.7,-0.8) rectangle ++(0.25,1.6);
    \node[below, font=\scriptsize, align=center] at (11.82,-1.75) {dense\\$128$};
    \filldraw[fill=green!35, draw=green!50!black] (13.0,-0.4) rectangle ++(0.25,0.8);
    \node[below, font=\scriptsize, align=center] at (13.12,-1.75) {logits\\$10$};
    \draw[->, thick] (9.7,0) -- (10.3,0);
    \draw[->, thick] (10.75,0) -- (11.6,0);
    \draw[->, thick] (12.05,0) -- (12.9,0);

    \draw[brace] (0,2.1) -- node[above=5pt, font=\small] {feature extractor} (9.6,2.1);
    \draw[brace] (10.3,2.1) -- node[above=5pt, font=\small] {classifier head} (13.4,2.1);
  \end{tikzpicture}
  \caption{A small CNN for MNIST. Each conv layer is followed by ReLU (not
    drawn). The feature maps get smaller and deeper; the final maps are
    flattened and classified by an MLP.}
  \label{fig:cnn}
\end{figure}

In PyTorch the whole network is a short \texttt{nn.Sequential}:

\begin{codebox}
model = nn.Sequential(
    nn.Conv2d(1, 16, kernel_size=3, padding=1),   # conv1
    nn.ReLU(),
    nn.MaxPool2d(2),                              # pool1
    nn.Conv2d(16, 32, kernel_size=3, padding=1),  # conv2
    nn.ReLU(),
    nn.MaxPool2d(2),                              # pool2
    nn.Flatten(),
    nn.Linear(32 * 7 * 7, 128),
    nn.ReLU(),
    nn.Linear(128, 10),                           # logits
)
\end{codebox}

As in Section 3.4, the model outputs logits; \texttt{nn.CrossEntropyLoss}
applies softmax and cross-entropy together.

\begin{stepbox}{Exercise 4 \quad Trace the network}
  Fill in the output shape (for one image) and the number of parameters of
  every layer in the model above.
  \begin{center}
    \renewcommand{\arraystretch}{1.35}
    \begin{tabular}{lll}
      \toprule
      Layer                              & Output shape                     & Parameters                                        \\
      \midrule
      Input                              & $(1, 28, 28)$                    & ---                                               \\
      Conv $1\!\to\!16$, $3\times3$, $p{=}1$ & $\sol[2.4cm]{(16, 28, 28)}$  & $\sol[5cm]{(1\cdot 9 + 1)\cdot 16 = 160}$         \\
      ReLU                               & $\sol[2.4cm]{(16, 28, 28)}$      & $\sol[5cm]{0}$                                    \\
      MaxPool $2\times2$                 & $\sol[2.4cm]{(16, 14, 14)}$      & $\sol[5cm]{0}$                                    \\
      Conv $16\!\to\!32$, $3\times3$, $p{=}1$ & $\sol[2.4cm]{(32, 14, 14)}$ & $\sol[5cm]{(16\cdot 9 + 1)\cdot 32 = 4\,640}$     \\
      MaxPool $2\times2$                 & $\sol[2.4cm]{(32, 7, 7)}$        & $\sol[5cm]{0}$                                    \\
      Flatten                            & $\sol[2.4cm]{(1568)}$            & $\sol[5cm]{0}$                                    \\
      Linear $\to 128$                   & $(128)$                          & $\sol[5cm]{(1568 + 1)\cdot 128 = 200\,832}$       \\
      Linear $\to 10$                    & $(10)$                           & $\sol[5cm]{(128 + 1)\cdot 10 = 1\,290}$           \\
      \midrule
      \multicolumn{2}{l}{Total}          & $\sol[5cm]{206\,922}$                                                                \\
      \bottomrule
    \end{tabular}
  \end{center}
  What percentage of the parameters is in the two convolutional layers?
  \[
    \sol[5cm]{4\,800 / 206\,922 \approx 2.3\%}
  \]
\end{stepbox}

Two observations:
\begin{itemize}
  \item This CNN has \emph{more} parameters than the MLP $784 \to 128 \to 10$
        (101\,770), yet it is more accurate on MNIST (around 99\% vs.\ 98\%)
        and much more robust to shifted digits. What matters is not only how
        many parameters a model has but how they are \emph{arranged}.
  \item Almost all the parameters sit in the first dense layer, which looks
        at the flattened $32 \times 7 \times 7$ maps. The convolutions, which do
        the actual feature extraction, are cheap.
\end{itemize}

% ──────────────────────────────────────────────────────────────────────────────
\subsection{Global Average Pooling}

The expensive dense layer can be removed. \textbf{Global average pooling}
(GAP) averages each channel over its \emph{entire} spatial extent, turning a
$(C, H, W)$ tensor into a vector of $C$ numbers, one per feature. A single
linear layer then maps these $C$ features to the logits. Most modern
architectures end this way. In PyTorch it is
\texttt{nn.AdaptiveAvgPool2d(1)} followed by \texttt{nn.Flatten()}.

\begin{stepbox}{Exercise 5 \quad Replace the dense head with GAP}
  \begin{enumerate}[label=(\alph*)]
    \item Apply GAP to a $(2, 2, 2)$ tensor with two channels:
          \[
            \left(
            \begin{pmatrix} 1 & 3 \\ 0 & 4 \end{pmatrix},\;
            \begin{pmatrix} 2 & 2 \\ 6 & 2 \end{pmatrix}
            \right)
            \;\xrightarrow{\ \text{GAP}\ }\;
            \bigl(\,\sol[1cm]{2},\ \sol[1cm]{3}\,\bigr)
          \]
    \item In the MNIST model, replace \texttt{Flatten} and both linear
          layers by GAP and one \texttt{nn.Linear(32, 10)}. The head now has
          $\sol[3cm]{(32 + 1) \cdot 10 = 330}$ parameters, and the whole
          network has $\sol[4cm]{160 + 4\,640 + 330 = 5\,130}$.
  \end{enumerate}
\end{stepbox}

% ══════════════════════════════════════════════════════════════════════════════
\section{Backpropagation Through Pooling}

Training a CNN means backpropagating through all three stages. Section 4.2
handled the convolution; ReLU passes the gradient where its input was
positive and blocks it elsewhere, as in Section 3.4. Pooling has no
parameters, but the gradient must still pass through it to reach the earlier
layers.
\begin{itemize}
  \item \textbf{Max pooling.} Only the largest value in each window affected
        the output, so the whole gradient goes to that position and the
        others in the window get $0$.
  \item \textbf{Average pooling.} Every value in a $k \times k$ window
        contributed $1/k^2$ of the output, so each receives $1/k^2$ of the
        gradient.
\end{itemize}

\begin{stepbox}{Exercise 6 \quad Route the gradient}
  The max-pooling output of Exercise 2 receives the gradient
  $\partial J / \partial(\text{output}) =
    \begin{pmatrix} 0.5 & -1 \\ 2 & 0.3 \end{pmatrix}$.
  Fill in $\partial J / \partial \bH$, using $\bH$ from Figure~\ref{fig:maxpool}.
  \[
    \pd{J}{\bH} =
    \begin{pmatrix}
      \sol[0.6cm]{0} & \sol[0.6cm]{0}   & \sol[0.6cm]{-1}  & \sol[0.6cm]{0} \\[2pt]
      \sol[0.6cm]{0} & \sol[0.6cm]{0.5} & \sol[0.6cm]{0}   & \sol[0.6cm]{0} \\[2pt]
      \sol[0.6cm]{0} & \sol[0.6cm]{0}   & \sol[0.6cm]{0.3} & \sol[0.6cm]{0} \\[2pt]
      \sol[0.6cm]{2} & \sol[0.6cm]{0}   & \sol[0.6cm]{0}   & \sol[0.6cm]{0}
    \end{pmatrix}
  \]
\end{stepbox}

% ══════════════════════════════════════════════════════════════════════════════
\section{What Does a CNN Learn?}

When we visualise what makes the units of a trained CNN respond, a clear
hierarchy appears (Zeiler and Fergus, 2014):
\begin{itemize}[nosep]
  \item \textbf{early layers} respond to edges, colour blobs and simple
        textures, much like the kernels of Section 4.2;
  \item \textbf{middle layers} respond to combinations of these: corners,
        circles, grids, fur, wheels;
  \item \textbf{late layers} respond to object parts and whole objects: eyes,
        faces, car fronts.
\end{itemize}
This is the hierarchy predicted in Section 4.1: each layer builds on the
features of the one below, over an ever larger receptive field.

\begin{notebox}
  \textbf{Try it.} The interactive \emph{CNN Explainer}
  (\url{https://poloclub.github.io/cnn-explainer/}) runs a small CNN in the
  browser and shows every convolution, ReLU and pooling step on real images.
\end{notebox}

% ══════════════════════════════════════════════════════════════════════════════
\section{Summary}

\begin{center}
  \renewcommand{\arraystretch}{1.35}
  \begin{tabular}{lll}
    \toprule
    Layer           & Parameters                                 & Effect on shape $(C, H, W)$                \\
    \midrule
    Conv2d          & $(C_{\text{in}} k^2 + 1)\, C_{\text{out}}$ & $C \to C_{\text{out}}$; $H, W$ by formula  \\
    ReLU            & 0                                          & unchanged                                  \\
    MaxPool2d (2)   & 0                                          & $H, W$ halved; $C$ unchanged               \\
    Global avg pool & 0                                          & $(C, H, W) \to (C)$                        \\
    Flatten         & 0                                          & $(C, H, W) \to (C H W)$                    \\
    Linear          & $(n_{\text{in}} + 1)\, n_{\text{out}}$     & $n_{\text{in}} \to n_{\text{out}}$         \\
    \bottomrule
  \end{tabular}
\end{center}

\begin{itemize}
  \item A CNN layer has three stages: convolution, ReLU, pooling.
  \item Max pooling keeps the largest value in each window. It has no
        parameters, works per channel, and makes the network approximately
        invariant to small shifts.
  \item A classification CNN is a feature extractor (conv blocks that shrink
        space and grow channels) followed by a classifier head (flatten or
        GAP, then linear layers, softmax and cross-entropy).
  \item In backpropagation, max pooling sends the gradient to the position of
        the maximum; average pooling spreads it evenly.
\end{itemize}

\paragraph{Coming up.}
Section 4.4 looks at the architectures that made CNNs famous, from LeNet-5
to ResNet, and at the design ideas each one introduced.

\end{document}
